% Independent D2 3D A-phi implementation, 2026-11-06.
% Not certified for infinite seawater or TEAM-7 reference accuracy.
\section{D2 三维 $A$--$\phi$ 直接法：边元、规范、边界和真实端口}\label{sec:bfzi-d2-aphi}

\paragraph{方法选择与独立性。}
承接第~\ref{sec:bfzi-physical-direct-locked}~节锁定的 BFZI 场景，
先用可独立检验的三维 Maxwell--MQS $A$--$\phi$ 有限元直接解，
不调用 G1 的五模态 Green/Coulomb 矩阵。
在 $\mu=\mu_0$、静止线性材料与 $e^{+\mathrm i\omega t}$ 基线，
固定两条真实闭合有向端口源曲线 $\Gamma_p$：
\begin{equation}
  \bm A_{\rm src,p}(\bm x)
    =\frac{\mu_0}{4\pi}\oint_{\Gamma_p}
        \frac{\mathrm d\bm\ell'}{|\bm x-\bm r'|},\quad
  \bm A=\sum_pI_p\bm A_{\rm src,p}+\bm a,\quad
  \bm E=-\mathrm i\omega\bm A-\nabla\phi.
\end{equation}
源中心线位于绝缘封装内部；薄丝只用于\emph{铜线外部}
海水耦合，不提供有限的铜自感、集肤或邻近效应。
圆形／圆角方形 BFZI 的单根连续螺旋和折线馈回流不得
被独立理想环集合或超椭圆线圈替代。

\subsection{反作用场与电荷守恒的混合弱式}
在由导电海水 $\Omega_c$ 和无电导绝缘/空气共同形成的有界
参考域 $Q_R$ 中，令 $\bm J_c=\sigma\bm E$：
\begin{align}
 \nabla\times\mu_0^{-1}\nabla\times\bm a
  &=\bm J_c,\qquad
 \nabla\cdot\bm J_c=0\quad(\Omega_c),\\
 \bm n\cdot\bm J_c&=0 \quad \text{于绝缘盒界面},\qquad
 \bm B_{\rm react}=\nabla\times\bm a .
\end{align}
令 $\bm v\in H(\mathrm{curl};Q_R)$、
$\psi\in H^1(\Omega_c)$，规范乘子 $p,q\in H^1_0(Q_R)$，
取一致的复共轭测试。反作用场的
$A$--$\phi$--$p$ 弱式为
\begin{align}
\int_{Q_R}\mu_0^{-1}(\nabla\times\bm a)\cdot
 (\nabla\times\bar{\bm v})\,\mathrm dV+
\int_{\Omega_c}\sigma(\mathrm i\omega\bm a+\nabla\phi)
 \cdot\bar{\bm v}\,\mathrm dV+
\int_{Q_R}\nabla p\cdot\bar{\bm v}\,\mathrm dV
 &=-\mathrm i\omega\int_{\Omega_c}\sigma\bm A_{\rm src}
 \cdot\bar{\bm v}\,\mathrm dV,\\
\int_{\Omega_c}\sigma
  \left(\bm a+\frac{\nabla\phi}{\mathrm i\omega}\right)
  \cdot\nabla\bar\psi\,\mathrm dV
 &=-\int_{\Omega_c}\sigma\bm A_{\rm src}
 \cdot\nabla\bar\psi\,\mathrm dV,\\
\int_{Q_R}\bm a\cdot\nabla\bar q\,\mathrm dV&=0 .
\end{align}
第二式是导电区零电荷累积约束除以
$\mathrm i\omega$，并自然给出导体--绝缘体界面
$\bm n\cdot\bm J_c=0$；它并不允许把任意材料界面
一律设为绝缘面。
$H(\mathrm{curl})$ 用最低阶有向 Nédélec 四面体边元，
$\phi$ 用导体上的连续 P1 标量元，规范乘子用
$H_0^1(Q_R)$ 标量元；对每个导体连通分量单独去除
$\phi$ 常数零模。该可运行初阶方法不能自动声称高精度；
需要材料共形网格、误差预算及 TEAM--7 测量比对。

\subsection{有限外边界不是无限海水：两种不同问题}
当前求解器具备两个**可选择的截断问题**：
\begin{align}
  \texttt{secondary\_zero:}&\quad \bm n\times\bm a=0,\\
  \texttt{total\_zero:}&\quad
     \bm n\times\left(\bm a+\bm A_{\rm src}\right)=0
  \quad\text{于 }\partial Q_R.
\end{align}
后者对全局边 $e=(a,b)$ 的自由度施加
\begin{equation}
  \int_e\bm a\cdot\mathrm d\bm\ell=
   -\int_e\bm A_{\rm src}\cdot\mathrm d\bm\ell ,
\end{equation}
使用真实有向边切向迹和独立的边上积分；
\emph{不是}在节点简单设置向量势零。
\texttt{total\_zero} 可用于与使用总矢势零边界的
Elmer TEAM--7 实例构造相同的有限域边界条件；
但即使如此，也仍须核对两软件的实体线圈源、
时间模型与外域尺寸。真实无限导电海水仍要求
至少三档外域（或透明边界/外场边界积分）证明
端口和空间场尾部的收敛。

\subsection{真实闭合端口、有限域边界项与功率}
对单位端口 $q$，反作用场 $\bm a_q$ 在源环上光滑，
真实海水\emph{增量感生电压}应从
\begin{equation}\label{eq:bfzi-d2-real-port-reaction}
 \boxed{\Delta Z_{pq}^{\mathrm{port}}
 =\mathrm i\omega\oint_{\Gamma_p}\bm a_q\cdot
 \mathrm d\bm\ell }
\end{equation}
抽取；该闭环积分对单值梯度规范变换不变。
多股/等效匝源需按其真实源电流方向和
安匝权重合成闭合路径，不能直接乘一个模糊的
等效半径。该感生电压与无限空间的
$\mathrm i\omega\int_{\Omega_c}
  \bm A_{\rm src,p}\cdot\bm J_{c,q}\,\mathrm dV$
在有限截断域中**不必相等**。
用源磁场 $\bm H_{\rm src,p}$、反作用磁场
$\bm H_{\rm react,q}$ 的连续积分分部关系，
两式的差含有边界项
\begin{equation}
 \begin{split}
  \oint_{\Gamma_p}\bm a_q\cdot\mathrm d\bm\ell
   - \int_{\Omega_c}\bm A_{\rm src,p}\cdot
      \bm J_{c,q}\,\mathrm dV
  =\int_{\partial Q_R}\big[
     (\bm n\times\bm H_{\rm src,p})\cdot\bm a_q
    -(\bm n\times\bm H_{\rm react,q})\cdot
       \bm A_{\rm src,p}
    \big]\,\mathrm dS .
 \end{split}
\end{equation}
此处源丝在域内、边界取不穿过源的光滑迹；
有限元离散还含线源曲线积分、材料网格和数值积分误差。
因此不能用共享同一体积求积的 Schur 补功率恒等式
宣称真实闭合端口精度已经得到独立验证。

海水的真实时间平均耗散始终是
\begin{equation}
 P_{\rm sea}(\bm I)=\frac12\int_{\Omega_c}\sigma
    \left|\sum_p I_p\bm E_p\right|^2\,\mathrm dV,
 \qquad q(\bm x;\bm I)=\frac12\sigma
    \left|\sum_p I_p\bm E_p(\bm x)\right|^2.
\end{equation}
在无限域正确衰减、全局源与材料分账一致、
边界通量已处理、数值误差趋零时，才可检验
$P_{\rm sea}=\frac12\bm I^{\mathrm H}
\Herm{\Delta\bm Z_{\rm sw}^{\mathrm{port}}}\bm I$。
粗网格或不同边界下二者有差距属于**未通过的数值误差预算**，
绝不可硬对称/截正端口矩阵去制造 PASS。

\paragraph{当前程序范围。}
\texttt{direct3d\_aphi.py} 解混合系统，
\texttt{direct3d\_port.py} 在各四面体内按边元插值
沿闭合物理端口做复电压积分；
\texttt{direct3d\_team7.py} 建立 COMPUMAG 偏心带孔
铝板及非 BFZI 的赛道形安匝源。
初步小网格求解和规范梯度极限已实跑，
但源路径反作用电压与体积 Joule 功率在粗网格
仍有不可忽略的差距。为独立方法对照，
可采用 ElmerCSC \texttt{CoilSolver/WhitneyAVSolver}
官方 TEAM--7 配置，必须匹配有限外域、总场边界、
$\pm2742$ A匝电流正方向和时域/频域观测定义。
\textbf{D2 尚无独立高保真 PASS；
G1 仍 FAIL，正式神经教师不得启动。}
